\section{Matériel végétal}
Trois cultures ont été étudiées : 

\begin{table}[H]
    \centering
    \caption{Variétés et modes de semis des trois cultures}
    \label{tab:varietes_semis}
    \begin{tabular}{|p{3cm}|p{5cm}|p{5.5cm}|}
        \hline
        \textbf{Culture} & \textbf{Variété} & \textbf{Mode de semis} \\
        \hline
        Laitue & EDEN & repiquage en motte \\
        \hline
        Navet & SUPER LONGO 2000 & Semis direct \\
        \hline
        Gombo & CLEMSON SPINELESS & Semis direct \\
        \hline
    \end{tabular}
\end{table}

\begin{figure}[H]
    \centering
    \imgstd{images/05_pepiniere_laitue.jpg}
    \caption{Pépinière de laitue variété EDEN. }
    \label{fig:laitue2}
\end{figure}

\begin{figure}[H]
    \centering
    \imgstd{images/06_semence_navet.jpg}
    \caption{Semence de Navet SUPER LONGO 2000. }
    \label{fig:navet2}
\end{figure}

\begin{figure}[H]
    \centering
    \imgstd{images/07_semence_gombo.jpg}
    \caption{Semence de Gombo CLEMSON SPINELESS. }
    \label{fig:gombo2}
\end{figure}

\subsection{Aménagement du sol et amendements}
Les amendements n'ont pas été pesés avec précision. Pour l'ensemble des trois plots TS, les quantités utilisées ont été les suivantes : 
un sac de compost conditionné dans un sac de riz de 50 kg, un demi-sac de biochar et de PCA, dont seule la moitié a été utilisée ; et 2 à 3 kg d'EGB. 
Le compost a été apporté en deux temps : une partie lors de l'aménagement des plots et une autre au cours de l'expérimentation.

